\section{Related Work}
\label{sec:related_work}

\subsection{MoE Inference Systems}
Early sparse-activation designs such as the Switch Transformer~\cite{fedus2021switch} demonstrated that routing each token to a small subset of experts substantially reduces computation while preserving model capacity. Subsequent serving systems strengthened this direction by addressing communication and scheduling costs at large scale~\cite{riquelme2021scaling, li2025uni, zhu2025exploring, zhang2025duoservemoe, yu2025prescope}. ExFlow~\cite{yao2024exploiting} leverages inter-layer expert affinity to colocate frequently co-activated experts, while TA-MoE~\cite{chen2022ta} uses device topology to reduce cross-node transfers. These approaches reduce communication but assume that expert availability can be largely predicted offline.

A separate line of work targets memory pressure on commodity or memory-constrained GPUs. ProMoE~\cite{song2024promoe}, ExpertFlow~\cite{he2024expertflow}, and Hobbit~\cite{tang2024hobbit} all employ cross-layer prefetching with a fixed forecasting horizon. ExpertFlow integrates token-aware expert prediction with proactive caching; Hobbit relies on mixed-precision offloading and multi-level caches; MoE-Lightning~\cite{cao2024moe} introduces hierarchical memory management to rebalance expert residency mid-execution. Other systems prioritize experts already in fast memory~\cite{bambhaniya2024moe} or inject null experts to smooth routing dynamics~\cite{zeng2024adamoe}. \systemname\ differs from all of the above by treating the prefetch horizon as a runtime control variable rather than an offline-tuned hyperparameter, by closing a feedback loop between stall/eviction signals and the horizon decision, and by letting cache-aware token routing absorb residual misprediction inside resident-expert compute.

\subsection{Routing and Load Balancing}
Routing quality is a central determinant of MoE efficiency~\cite{huang2025dynamic, zhou2025moe}. Expert Choice~\cite{zhou2022mixture} provides load balancing by letting experts self-select tokens; AdaMOE~\cite{zeng2024adamoe} dynamically adjusts the number of active experts per token; ExFlow~\cite{he2024expertflow} maintains routing consistency across layers to reduce communication; MoDSE~\cite{sun2024mixture} assigns heterogeneous expert sizes to match input complexity. These techniques shape \emph{which} experts are activated. \systemname\ is orthogonal: rather than designing new routing objectives or expert topologies, we treat the cross-layer prediction interval as a runtime control variable, leaving routing semantics unchanged. The two directions are complementary, and \systemname\ can be deployed in front of any of the routing schemes above.
